\section{Limitaciones metodológicas}

Hay tres advertencias que condicionan cómo deben leerse los números de este informe.

La primera es que el pico de RAM no resulta comparable entre macOS y Windows. Pandas en XL reporta 4,522 MB en Pc\_Jesus y 12,096 MB en Pc\_Nahin para el mismo código y los mismos datos. macOS comprime páginas de memoria y el RSS que reporta \texttt{psutil} no las contabiliza, así que estas cifras sirven para comparar herramientas \textbf{dentro} de una misma máquina, pero no entre sistemas operativos.

La segunda es que el equipo de mayores especificaciones resultó el más lento. Pc\_Nahin (Core Ultra 9, 40 GB) tarda de 2 a 4 veces más que Pc\_Jesus (8 GB) en todas las herramientas, incluida la carga pura de CSV con DuckDB, que pasa de 8 a 36 segundos en XL. La brecha es de entrada/salida y no de CPU ni de memoria. Las causas plausibles son una carpeta sincronizada con la nube, un antivirus escaneando 4.7 GB en cada lectura, o las tres versiones corridas en una sola sesión sin liberar memoria. No se verificó cuál aplica, así que los tiempos absolutos de ese equipo deben tomarse como cota superior.

La tercera es que las versiones de librerías difieren entre equipos (Cuadro \ref{tab:software}). Eso no explica diferencias de orden de magnitud, pero sí introduce ruido en las comparaciones finas. Por eso las conclusiones de este informe se apoyan en órdenes de magnitud y en fallos categóricos, y no en márgenes estrechos.
